\documentclass[10pt,a4paper]{amsart}
\usepackage[margin=19mm]{geometry}
\usepackage{amsmath,amssymb,mathtools}
\usepackage[scr=rsfs,cal=euler]{mathalfa}
\usepackage{xfrac}
\usepackage[unicode,hidelinks,bookmarks=false]{hyperref}
\newcommand{\ie}{i.e.\ }
\newcommand{\cf}{cf.\ }
\newcommand{\resp}{respectively\ }
\newcommand{\Subsection}{\subsection}
\providecommand{\nobreakdash}{\nobreak\hbox{-}}
\setlength{\emergencystretch}{2em}
\multlinegap0pt
\usepackage{bm}
\usepackage{bbm}
\usepackage{dsfont}
\usepackage{xspace}
\usepackage{cancel}
\usepackage{mathtools}
\usepackage[shortlabels]{enumitem}
\usepackage{amsthm}
\makeatletter
\@ifundefined{theorem}{\newtheorem{theorem}{Theorem}}{}
\@ifundefined{proposition}{\newtheorem{proposition}[theorem]{Proposition}}{}
\makeatother
\DeclareMathOperator{\proj}{proj}
\title[Carath\'eodory number of homogeneous convex cones]{Carath\'eodory number of homogeneous convex cones}
\author{Chek Beng Chua}
\date{Source: arXiv:2511.17051v2 (CC BY 4.0)}
\begin{document}
\maketitle

\noindent\emph{Source and licence.} Carath\'eodory number of homogeneous convex cones. Version arXiv:2511.17051v2 (CC BY 4.0), distributed under the Creative Commons Attribution 4.0 International licence, as recorded in the arXiv licence metadata for this exact version and shown on the abstract page https://arxiv.org/abs/2511.17051v2. Full rights notice: licenses/arxiv-2511.17051-CC-BY-4.0.txt.\par\smallskip

\noindent\textit{Editorial scope.} Counted unit: the Theorem whose source label is \texttt{theorem:characterizationofrankmatchingcaratheodorynumber} in the article (source0.tex lines 1682--1699, source label \texttt{theorem:characterizationofrankmatchingcaratheodorynumber}), delivered together with the article's own complete original proof of it (source0.tex lines 1700--1839, one \texttt{proof} environment of 140 lines). No mathematics is written, repaired or re-derived: statement and proof are the article's own text line for line. Required context retained uncounted: identity:leftadjointleft (lines 1092--1124); identity:leftadjointself (lines 1151--1162); proposition:Caratheodorynumbersnecessary (lines 1389--1429); proposition:Caratheodorynumberssufficient (lines 1540--1561). Every definition, display and label of those lines is retained unchanged. Omitted from this package and not redistributed: 1--1091, 1125--1150, 1163--1388, 1430--1539, 1562--1681, 1840--2756. Nothing inside a retained excerpt is omitted or altered: every retained body is delivered exactly as the article's lines are written, so no omission receipt of any kind accompanies this package. Citations: the retained excerpts carry no citation command at all, so no bibliography entry of the article is transcribed and no reference is redistributed. Reference adaptation: none was needed and none was made; every reference inside the delivered unit resolves inside it, so no unresolved reference or citation is produced. Printed slips kept literal: no printed word, symbol, hypothesis or number of the counted statement or of its proof was altered. Layout and class: the article's own class is \texttt{amsart} with packages including amssymb, geometry, enumitem, hyperref; the wrapper is the lane's pinned \texttt{amsart} editorial wrapper at 10pt on A4, which restates the theorem-like environments the retained text needs and adds the article's own macro definitions that the retained text needs (\texttt{\textbackslash proj}), with extra wrapper packages (bm, bbm, dsfont, xspace, cancel, mathtools, enumitem). The delivered numbering is the unit's own. Assets: no retained span contains a figure environment, a graphics-inclusion command, a PDF-page inclusion, a table or a TikZ picture; no image file is copied into this package, the package declares no asset dependency and no image file is redistributed with it. Licence: version arXiv:2511.17051v2 of this article is distributed under the Creative Commons Attribution 4.0 International licence, as recorded in the arXiv licence metadata for this exact version and shown on the abstract page https://arxiv.org/abs/2511.17051v2. A full rights notice is delivered at licenses/arxiv-2511.17051-CC-BY-4.0.txt. Characters: every retained excerpt is pure ASCII, so the delivered engine, XeLaTeX with its default Latin Modern text font, reports no missing character.
\begin{enumerate}[label={[\Alph*]},
ref={[\Alph*]}]
\item\label{identity:leftleft}%
$\mathcal{L}_{A_{ij}} \mathcal{L}_{B_{jk}} =
\mathcal{L}_{\mathcal{L}_{A_{ij}} B_{jk}}$ and, by taking the adjoint,
$\mathcal{L}^{*}_{B_{jk}} \mathcal{L}^{*}_{A_{ij}} =
\mathcal{L}^{*}_{\mathcal{L}_{A_{ij}} B_{jk}}$.

\item\label{identity:leftright}%
$\mathcal{L}_{A_{ij}} \mathcal{R}_{C_{kl}} = \mathcal{R}_{C_{kl}}
\mathcal{L}_{A_{ij}}$ and, by taking the adjoint, $\mathcal{R}^{*}_{C_{kl}}
\mathcal{L}^{*}_{A_{ij}} = \mathcal{L}^{*}_{A_{ij}}
\mathcal{R}^{*}_{C_{kl}}$.

\item\label{identity:rightright}%
$\mathcal{R}_{C_{kl}} \mathcal{R}_{B_{jk}} =
\mathcal{R}_{\mathcal{R}_{C_{kl}} B_{jk}}$ and, by taking the adjoint,
$\mathcal{R}^{*}_{B_{jk}} \mathcal{R}^{*}_{C_{kl}} =
\mathcal{R}^{*}_{\mathcal{R}_{C_{kl}} B_{jk}}$.

\item\label{identity:leftadjointleft}%
$\mathcal{L}^{*}_{A_{ij}} \mathcal{L}_{D_{ik}} =
\mathcal{L}_{\mathcal{L}^{*}_{A_{ij}} D_{ik}}$ if $j < k$ and, by taking
the adjoint, $\mathcal{L}^{*}_{D_{ik}} \mathcal{L}_{A_{ij}} =
\mathcal{L}^{*}_{\mathcal{L}^{*}_{A_{ij}} D_{ik}}$.

\item\label{identity:leftadjointright}%
$\mathcal{L}^{*}_{A_{ij}} \mathcal{R}_{C_{kl}} = \mathcal{R}_{C_{kl}}
\mathcal{L}^{*}_{A_{ij}}$ if $j < k$ and, by taking the adjoint,
$\mathcal{R}^{*}_{C_{kl}} \mathcal{L}_{A_{ij}} = \mathcal{L}_{A_{ij}}
\mathcal{R}^{*}_{C_{kl}}$.

\end{enumerate}

\begin{enumerate}[label={[\Alph*]},
ref={[\Alph*]}]
\addtocounter{enumi}{6}%
\item\label{identity:leftadjointself}%
$\mathcal{L}^{*}_{A_{ij}} \mathcal{L}_{A_{ij}} = \frac{1}{n_{j}}
\left\lVert A_{ij}\right\rVert_{F}^{2} \mathcal{I}_{\mathbb{V}_{jk}}$ and,
by polarization, $\mathcal{L}^{*}_{A_{ij}} \mathcal{L}_{B_{ij}} +
\mathcal{L}^{*}_{B_{ij}} \mathcal{L}_{A_{ij}} = \frac{2}{n_{j}}
\left\langle A_{ij} \middle| B_{ij} \right\rangle_{\mathrm{F}}
\mathcal{I}_{\mathbb{V}_{jk}}$,

\end{enumerate}

\begin{proposition}\label{proposition:Caratheodorynumbersnecessary}
For a homogeneous convex cone of rank $r$, if the Carath\'eodory number of
its closure is $r$, then, in any Ishi spectrahedral representation
$\mathbb{S}^{n}_{++} \cap \mathbb{V}$ of the homogeneous cone, for all $1
\leq i < j < k \leq r$, either of the following two equivalent conditions
hold:
\begin{enumerate}[label=(\roman*)]
\item For all $A_{ij} \in \mathbb{V}_{ij}$,
\[
\left\lVert A_{ij}\right\rVert_{F}^{2} \left\lVert
B_{ik}\right\rVert_{F}^{2} = n_{j} \left\lVert \mathcal{L}^{*}_{A_{ij}}
B_{ik}\right\rVert_{F}^{2} \quad \forall B_{ik} \in \mathbb{V}_{ik},
\]
or equivalently, $\mathcal{L}_{A_{ij}} \mathcal{L}^{*}_{A_{ij}} =
\frac{1}{n_{j}} \left\lVert A_{ij}\right\rVert_{F}^{2}
\mathcal{I}_{\mathbb{V}_{ik}}$.

\item $\mathbb{V}_{ik}$ and $\mathbb{V}_{jk}$ have the same dimensions
whenever $\mathbb{V}_{ij}$ is not the trivial subspace.

\end{enumerate}
Dually, if the Carath\'eodory number of its dual cone is $r$, then, in any
Ishi spectrahedral representation $\mathbb{S}^{n}_{++} \cap \mathbb{V}$ of
the homogeneous cone, for all $1 \leq k < j < i \leq r$, either of the
following two equivalent conditions hold:
\begin{enumerate}[label=(\roman*)]
\item For all $A_{ji} \in \mathbb{V}_{ji}$,
\[
\left\lVert A_{ji}\right\rVert_{F}^{2} \left\lVert
B_{ki}\right\rVert_{F}^{2} = n_{j} \left\lVert \mathcal{R}^{*}_{A_{ji}}
B_{ki}\right\rVert_{F}^{2} \quad \forall B_{ki} \in \mathbb{V}_{ki},
\]
or equivalently, $\mathcal{R}_{A_{ji}} \mathcal{R}^{*}_{A_{ji}} =
\frac{1}{n_{j}} \left\lVert A_{ji}\right\rVert_{F}^{2}
\mathcal{I}_{\mathbb{V}_{ki}}$.

\item $\mathbb{V}_{ki}$ and $\mathbb{V}_{kj}$ have the same dimensions
whenever $\mathbb{V}_{ji}$ is not the trivial subspace.

\end{enumerate}
\end{proposition}

\begin{proposition}\label{proposition:Caratheodorynumberssufficient}
For a homogeneous convex cone of rank $r$, if, in any Ishi spectrahedral
representation $\mathbb{S}^{n}_{++} \cap \mathbb{V}$ of the homogeneous
cone, for every $i \in \left\{1, \dotsc, r\right\}$, the orthogonal
projection of every extreme ray of the form
$\mathsf{F}_{T,\left\{i\right\}}$ onto the linear span of every face of the
form $\mathsf{F}_{I_{n},\mathsf{B}}$ with $\mathsf{B} \subseteq \left\{i+1,
\dotsc, r\right\}$ lies inside an extreme ray of the closure
$\mathbb{S}^{n}_{+} \cap \mathbb{V}$, then the Carath\'eodory number of
each nontrivial face of the closure matches the rank of its homogeneous
relative interior.

Dually, if, in any Ishi spectrahedral representation $\mathbb{S}^{n}_{++}
\cap \mathbb{V}$ of the homogeneous cone, for every $i \in \left\{1,
\dotsc, r\right\}$, the orthogonal projection of every extreme ray of the
form $\mathsf{F}^{*}_{T,\left\{i\right\}}$ onto the linear span of every
face of the form $\mathsf{F}^{*}_{I_{n},\mathsf{N}}$ with $\mathsf{N}
\subseteq \left\{1, \dotsc, i-1\right\}$ lies inside an extreme ray of the
dual cone $\proj_{\mathbb{V}} \mathbb{S}^{n}_{+}$, then the Carath\'eodory
number of each nontrivial face of the dual cone matches the rank of its
homogeneous relative interior.
\end{proposition}

\begin{theorem}\label{theorem:characterizationofrankmatchingcaratheodorynumber}
For every homogeneous convex cone of rank $r$, the Carath\'eodory number of
its closure is $r$ if and only if, in any Ishi spectrahedral representation
$\mathsf{K} = \mathbb{S}^{n}_{++} \cap \mathbb{V}$ of the homogeneous
convex cone, for all $1 \leq i < j < k \leq r$, $\mathbb{V}_{ik}$ and
$\mathbb{V}_{jk}$ have the same dimensions whenever $\mathbb{V}_{ij}$ is
not the trivial subspace. In this case, the Carath\'eodory number of every
nontrivial face of its closure matches the rank of its homogeneous relative
interior.

Dually, the Carath\'eodory number of its dual cone is $r$ if and only if,
in any Ishi spectrahedral representation $\mathsf{K} = \mathbb{S}^{n}_{++}
\cap \mathbb{V}$ of the homogeneous convex cone, for all $1 \leq k < j < i
\leq r$, $\mathbb{V}_{ki}$ and $\mathbb{V}_{kj}$ have the same dimensions
whenever $\mathbb{V}_{ji}$ is not the trivial subspace. In this case, the
Carath\'eodory number of every nontrivial face of its dual cone matches the
rank of its homogeneous relative interior.
\end{theorem}

\begin{proof}
We have previously established in
Proposition~\ref{proposition:Caratheodorynumbersnecessary} the equivalence
of the respective dimension characterizations of rank-matching
Carath\'eodory numbers in the theorem statement to
\[
\forall 1 \leq i < j < k \leq r \ \forall A_{ij} \in \mathbb{V}_{ij} \quad
\mathcal{L}_{A_{ij}} \mathcal{L}^{*}_{A_{ij}} = \frac{1}{n_{j}} \left\lVert
A_{ij}\right\rVert_{F}^{2} \mathcal{I}_{\mathbb{V}_{ik}}
\]
and
\[
\forall 1 \leq k < j < i \leq r \ \forall A_{ji} \in \mathbb{V}_{ji} \quad
\mathcal{R}_{A_{ji}} \mathcal{R}^{*}_{A_{ji}} = \frac{1}{n_{j}} \left\lVert
A_{ji}\right\rVert_{F}^{2} \mathcal{I}_{\mathbb{V}_{ki}}.
\]
With the proofs of necessity in
Proposition~\ref{proposition:Caratheodorynumbersnecessary} and sufficiency
in Proposition~\ref{proposition:Caratheodorynumberssufficient}, it remains
to show that
\begin{enumerate}[label=(\roman*)]
\item for any $i \in \left\{1, \dotsc, r-1\right\}$, if
$\mathcal{L}_{A_{ij}} \mathcal{L}^{*}_{A_{ij}} = \frac{1}{n_{j}}
\left\lVert A_{ij}\right\rVert_{F}^{2} \mathcal{I}_{\mathbb{V}_{ik}}$ for
all $A_{ij} \in \mathbb{V}_{ij}$ with $j \in \left\{i+1, \dotsc, r\right\}$
and all $k \in \left\{j+1, \dotsc, r\right\}$, then the orthogonal
projection of any extreme ray of the form $\mathsf{F}_{T,\left\{i\right\}}$
onto the linear span of the face $\mathsf{F}_{I_{n},\mathsf{B}}$ with
$\mathsf{B} \subseteq \left\{i+1, \dotsc, r\right\}$ lies in an extreme
ray; and

\item for any $i \in \left\{2, \dotsc, r\right\}$, if $\mathcal{R}_{A_{ji}}
\mathcal{R}^{*}_{A_{ji}} = \frac{1}{n_{j}} \left\lVert
A_{ji}\right\rVert_{F}^{2} \mathcal{I}_{\mathbb{V}_{ki}}$ for all $A_{ji}
\in \mathbb{V}_{ji}$ with $j \in \left\{1, \dotsc, i-1\right\}$ and all $k
\in \left\{1, \dotsc, j-1\right\}$, then the orthogonal projection of any
extreme ray of the form $\mathsf{F}^{*}_{T,\left\{i\right\}}$ onto the
linear span of the face $\mathsf{F}^{*}_{I_{n},\mathsf{N}}$ with
$\mathsf{N} \subseteq \left\{1, \dotsc, i-1\right\}$ lies in an extreme
ray.

\end{enumerate}

To prove the first statement, we suppose that $\mathcal{L}_{A_{ij}}
\mathcal{L}^{*}_{A_{ij}} = \frac{1}{n_{j}} \left\lVert
A_{ij}\right\rVert_{F}^{2} \mathcal{I}_{\mathbb{V}_{ik}}$ for all $A_{ij}
\in \mathbb{V}_{ij}$ with $j \in \left\{i+1, \dotsc, r\right\}$ and all $k
\in \left\{j+1, \dotsc, r\right\}$ and that $\mathsf{B}$ is a nonempty
subset of $\left\{i+1, \dotsc, r\right\}$, and consider an extreme ray
$\mathsf{F}_{T,\left\{i\right\}} = \mathbb{R}_+(T^{\top}
I^{\left\{i\right\}} T)$. If the blocks $T_{ij}$ for $j \in \mathsf{B}$ are
all zero blocks, then $T^{\top} I^{\left\{i\right\}} T =
I^{\left\{i\right\}}$, and hence its orthogonal projection onto the linear
span of the face $\mathsf{F}_{I_{n},\mathsf{B}} = \mathbb{S}^{n}_{+} \cap
\mathbb{V}^{\mathsf{B}}$ is in the trivial face, which is in every extreme
ray.

Otherwise, we let $j \in \mathsf{B}$ denote the least index for which
$T_{ij}$ is a nonzero block. In this case, the orthogonal projection of
$T^{\top} I^{\left\{i\right\}} T$ is $\hat{T}^{\top} \hat{T}$, where
$\hat{T}$ has $\hat{T}_{ik} = T_{ik}$ for $k \in \left\{j, \dotsc,
r\right\} \cap \mathsf{B}$ and zero blocks everywhere else. We now show
that $\hat{T}^{\top} \hat{T} = \mathcal{T}_{U} I^{\left\{j\right\}} =
U^{\top} I^{\left\{j\right\}} U$, where $\mathcal{T}_{U} \in
\mathsf{T}_{\mathsf{K}}$ has $U_{jj} = \sqrt{n_{j}}^{-1} \left\lVert
T_{ij}\right\rVert_{F} I_{n_{j}}$, $U_{jk} = \sqrt{n_{j}} \left\lVert
T_{ij}\right\rVert_{F}^{-1} \mathcal{L}^{*}_{T_{ij}} T_{ik}$ for $k \in
\left\{j+1, \dotsc, r\right\} \cap \mathsf{B}$, and zero off-diagonal
blocks everywhere else. Indeed,
\begin{enumerate}[label=(\roman*)]
\item in the rows not indexed by $\left\{j, \dotsc, r\right\} \cap
\mathsf{B}$, both $\hat{T}^{\top} \hat{T}$ and $U^{\top}
I^{\left\{j\right\}} U$ have zero blocks;

\item the $j$-th diagonal blocks of $\hat{T}^{\top} \hat{T}$ and $U^{\top}
I^{\left\{j\right\}} U$ agree because
\[
\mathcal{L}^{*}_{U_{jj}} \mathcal{L}_{U_{jj}} I_{n_{j}} =
\frac{1}{n_{j}}\left\lVert T_{ij}\right\rVert_{F}^{2} I_{n_{j}}
\overset{\ref{identity:leftadjointself}}{=} \mathcal{L}^{*}_{T_{ij}}
\mathcal{L}_{T_{ij}} I_{n_{j}};
\]

\item for each $k \in \left\{j+1, \dotsc, r\right\} \cap \mathsf{B}$, the
$(j, k)$-th blocks of $\hat{T}^{\top} \hat{T}$ and $U^{\top}
I^{\left\{j\right\}} U$ agree because
\[
\mathcal{L}^{*}_{U_{jj}} U_{jk} = \sqrt{n_{j}}^{-1} \left\lVert
T_{ij}\right\rVert_{F} \sqrt{n_{j}} \left\lVert T_{ij}\right\rVert_{F}^{-1}
\mathcal{L}^{*}_{I_{n_{j}}} \mathcal{L}^{*}_{T_{ij}} T_{ik} =
\mathcal{L}^{*}_{T_{ij}} T_{ik};
\]
and, finally,

\item for any $k \in \left\{j+1, \dotsc, r\right\} \cap \mathsf{B}$ and any
$l \in \left\{j+1, \dotsc, k\right\} \cap \mathsf{B}$, the $(l, k)$-th
blocks of $\hat{T}^{\top} \hat{T}$ and $U^{\top} I^{\left\{j\right\}} U$
agree because
\[
\begin{split}
\mathcal{L}^{*}_{U_{jl}} U_{jk}
%
&= n_{j} \left\lVert T_{ij}\right\rVert_{F}^{-2}
\mathcal{L}^{*}_{\mathcal{L}^{*}_{T_{ij}} T_{il}} \mathcal{L}^{*}_{T_{ij}}
T_{ik}
\\
&\overset{\ref{identity:leftadjointleft}}{=} n_{j} \left\lVert
T_{ij}\right\rVert_{F}^{-2} \mathcal{L}^{*}_{T_{il}} \mathcal{L}_{T_{ij}}
\mathcal{L}^{*}_{T_{ij}} T_{ik}
\\
&= \mathcal{L}^{*}_{T_{il}} \mathcal{I}_{\mathbb{V}_{ik}} T_{ik}
%
= \mathcal{L}^{*}_{T_{il}} T_{ik},
\end{split}
\]
where the third equality follows from the hypothesis $\mathcal{L}_{A_{ij}}
\mathcal{L}^{*}_{A_{ij}} = \frac{1}{n_{j}} \left\lVert
A_{ij}\right\rVert_{F}^{2} \mathcal{I}_{\mathbb{V}_{ik}}$.

\end{enumerate}
Therefore, the orthogonal projection of $T^{\top} I^{\left\{i\right\}} T$,
and hence of the face $\mathsf{F}_{T,\left\{i\right\}}$, onto the linear
span of the face $\mathsf{F}_{I_{n},\mathsf{B}}$ is the extreme ray
$\mathsf{F}_{U,\left\{j\right\}}$.

The proof of the statement for the dual cone is similar, where, in the
nontrivial case, the orthogonal projection of $\proj_{\mathbb{V}}\left( T
I^{\left\{i\right\}} T^{\top} \right)$ is $\proj_{\mathbb{V}}\left( \hat{T}
\hat{T}^{\top} \right)$ with $\hat{T}$ having $\hat{T}_{ki} = T_{ki}$ for
$k \in \left\{1, \dotsc, j\right\} \cap \mathsf{N}$ and zero blocks
everywhere else; which can then be shown to be the same as
$\mathcal{T}^{*}_{U} I^{\left\{j\right\}} = \proj_{\mathbb{V}}\left( U
I^{\left\{j\right\}} U^{\top} \right)$, with $U \in \mathbb{T}$ having
$U_{jj} = \sqrt{n_{j}}^{-1} \left\lVert T_{ji}\right\rVert_{F} I_{n_{j}}$,
$U_{kj} = \sqrt{n_{j}} \left\lVert T_{ji}\right\rVert_{F}^{-1}
\mathcal{R}^{*}_{T_{ji}} T_{ki}$ for $k \in \left\{1, \dotsc, j-1\right\}
\cap \mathsf{N}$, and zero off-diagonal blocks everywhere else, using a
similar argument with the left multiplication operators replaced by
corresponding right multiplication operators.
\end{proof}

\end{document}
